%% kapa: exam=1 round=36 year=2025 period=2 form=A subject=appraisal_law subject_ko=감정평가관계법규 seq=18 number=18
%% source: Q-Net 감정평가사 기출문제 (cst00309) artlSeq=5231525 file=raw/1008/5231525/2025년 제36회 감정평가사 1차시험 2교시 문제지 원본.pdf
%% conversion: kordoc 4.24.51 -> tools/gen1.py (auto); review=proofread-sonnet; flags=none
\begin{question}{18}
감정평가 및 감정평가사에 관한 법령상 감정평가에 관한 설명으로 옳지 않은 것은?
\begin{choices}
  \item 국토교통부장관은 타당성조사의 대상이 되는 감정평가를 한 감정평가사에 대하여 징계처분을 할 수 없어 타당성조사의 실익이 없는 경우에는 조사를 하지 않을 수 있다.
  \item 감정평가법인은 감정평가서를 의뢰인에게 발급하기 전에 감정평가를 한 소속 감정평가사가 작성한 감정평가서의 적정성을 같은 법인 소속의 다른 감정평가사에게 심사하게 하여야 한다.
  \item 국가는 토지등의 관리를 위하여 토지등을 감정평가하려는 경우에는 감정평가법인등에 의뢰하여야 한다.
  \item 신탁회사가 자산의 매입ㆍ매각ㆍ관리를 위하여 감정평가를 의뢰하려는 경우 한국감정평가사협회에 요청하여 추천받은 감정평가법인등에 감정평가를 의뢰할 수 있다.
  \item 감정평가법인등은 감정평가서의 원본을 발급일부터 2년 이상 보존하여야 한다.
\end{choices}
\end{question}
\answer{5}
